%%%%%%%%%%%%%%%%%%%%%%%%%%%%%%%%%%%%%%%%%%%%%%%%%%%%%%%%%%%%%%
%% Chapter 5: Text Translation (gloss2text / text2gloss)
%%%%%%%%%%%%%%%%%%%%%%%%%%%%%%%%%%%%%%%%%%%%%%%%%%%%%%%%%%%%%%

\section{Text Translation (gloss2text / text2gloss)}
\label{chap:text-translation}

Chapter~\ref{chap:isolated-gloss-recognition} addressed the video-to-gloss
stage of the pipeline introduced in
Section~\ref{sec:motivation-context}: converting a segmented sign video into a
gloss label. This chapter addresses the remaining symbolic-to-natural-language
step, translating between gloss sequences and natural-language text in both
directions -- gloss-to-text and text-to-gloss -- so that the recognised gloss
sequences can be rendered as fluent sentences, and so that the reverse direction
can be used for data augmentation and evaluation.

\todo{PLACEHOLDER: This chapter is not yet covered by the project's existing
research and development reports and needs to be written once the
gloss-text translation experiments are conducted. The subsections below outline
the expected structure.}

\subsection{Approach and data}
\label{sec:text-translation-approach-data}

\todo{PLACEHOLDER: Describe the parallel gloss-text corpus (or corpora) used for
training and evaluation, how gloss sequences are derived from the
video-to-gloss output of Chapter~\ref{chap:isolated-gloss-recognition} or from
existing annotated corpora, any data cleaning/normalisation applied to gloss and
text sequences, and the train/validation/test split methodology.}

\subsection{Model and training}
\label{sec:text-translation-model-training}

\todo{PLACEHOLDER: Describe the sequence-to-sequence architecture used for
gloss2text and text2gloss translation (e.g.\ Transformer encoder-decoder,
pretrained language model fine-tuning), tokenisation strategy for gloss tokens,
training objective, optimiser, and hyperparameters. Relate architectural choices
back to the NLP models surveyed in Section~\ref{sec:nlp-models-sign-language}.}

\subsection{Evaluation (both directions)}
\label{sec:text-translation-evaluation}

\todo{PLACEHOLDER: Report quantitative results for both gloss-to-text and
text-to-gloss directions (e.g.\ BLEU, ROUGE, or other machine-translation
metrics as appropriate), qualitative examples of translated output, and an
analysis of failure modes. This evaluation will be reused as the text
translation stage of the end-to-end pipeline evaluated in
Chapter~\ref{chap:sentence-level-translation}.}
